\documentclass[11pt]{article}

\usepackage[utf8]{inputenc}
\usepackage[T1]{fontenc}
\usepackage{amsmath,amssymb}
\usepackage{graphicx}
\usepackage{booktabs}
\usepackage{siunitx}
\usepackage[hidelinks]{hyperref}
\usepackage[round]{natbib}
\usepackage{geometry}
\geometry{margin=1in}

\graphicspath{{figures/}}
% Compile-safe figure include: shows a placeholder box if the image is absent,
% so the paper builds before the figures are copied into paper/figures/.
\newcommand{\safefig}[2][width=0.8\linewidth]{%
  \IfFileExists{figures/#2}{\includegraphics[#1]{#2}}%
  {\fbox{\parbox[c][3cm][c]{0.8\linewidth}{\centering\ttfamily figure placeholder: figures/#2}}}}

\title{Can Explainability Detect Shortcut Learning?\\
\large A Controlled Dose--Response Study of Grad-CAM, Integrated Gradients,\\ and Occlusion Sensitivity on a Medical Image Classifier}
\author{Aanya Mittal}
\date{\today}

\begin{document}
\maketitle

\begin{abstract}
% Grounded in the actual experiment. Revise wording freely; keep the numbers.
Deep networks for medical imaging often rely on \emph{shortcuts}---spurious cues correlated
with the label---rather than disease-relevant features. We ask whether three widely used
post-hoc attribution methods (Grad-CAM, Integrated Gradients, and Occlusion Sensitivity) can
detect such reliance under controlled conditions. We inject a small artificial marker at a
random location into PneumoniaMNIST images, correlating it with the disease label at four
strengths ($\rho \in \{0.0, 0.7, 0.9, 1.0\}$), train twelve classifiers (four strengths,
three seeds), and verify behavioral reliance before trusting any explanation. We measure
attribution mass above a random-box chance level ($\mathrm{AM}^{\star}$), pointing-game hit
rate, and a detection rule that requires both above-chance localization and measured behavioral
reliance. Integrated Gradients and Occlusion Sensitivity localize the marker in all nine
shortcut-reliant models and in none of the three controls, with $\mathrm{AM}^{\star}$ rising
monotonically with $\rho$ (Spearman $+0.95$ and $+0.89$ against $\rho$). Grad-CAM fails to
localize the marker in all twelve models; however, we show this outcome is confounded by the
coarse spatial resolution of its feature map on our architecture and cannot be read as evidence
that Grad-CAM is blind to shortcuts. We also observe that Integrated Gradients places
above-chance mass on the marker even in control models that do not rely on it, underscoring that
attribution alone is not evidence of reliance. These are descriptive findings from three seeds
on a single binary task and are not claimed to be statistically powered or to generalize.
\end{abstract}

% ===========================================================================
\section{Introduction}
\label{sec:intro}

Deep neural networks often reach high benchmark accuracy by exploiting \emph{shortcuts}: decision
rules that succeed on the training distribution but fail to transfer to more challenging conditions
\citep{geirhos2020shortcut}. Rather than learning the intended concept, a model may latch onto an
incidental cue that happens to correlate with the label. This matters most where errors are costly,
and medical imaging is a clear example. Classifiers trained to read chest radiographs have been
shown to infer the source hospital---and with it disease prevalence---rather than pathology
\citep{zech2018variable}; COVID-19 detectors have been found to rely on acquisition and dataset-source
confounders instead of lung findings \citep{degrave2021ai}; and a melanoma classifier was swayed by
surgical skin markings a clinician had drawn near a benign lesion \citep{winkler2019association}.
In each case the model was accurate on data resembling its training set and unreliable beyond it.

Because such failures are invisible to ordinary accuracy metrics, feature-attribution methods from
explainable AI (XAI) are frequently proposed as a safeguard: if an attribution map highlights an
irrelevant region, the reasoning goes, a practitioner can catch the shortcut before deployment.
Whether attributions actually support this use is contested. Several studies report that popular
saliency methods fail basic sanity checks \citep{adebayo2018sanity}, localize pathology poorly in
medical images \citep{arun2021assessing,saporta2022benchmarking}, and---most directly---are
unreliable for surfacing injected spurious signals, particularly when those signals are unknown or
not visually obvious \citep{adebayo2020debugging,adebayo2022posthoc}. The practical question is
therefore not whether attributions \emph{can} be computed, but whether, and under what conditions,
they \emph{faithfully reveal} a model's reliance on a shortcut.

We study this question in a deliberately controlled setting. We inject an artificial marker into a
binary chest-X-ray task and correlate it with the label at several strengths, so that the shortcut
is known by construction and detection can be scored against ground truth. We stress at the outset
that the broad question ``can explainability detect shortcut learning?'' is \emph{not} novel: it is
essentially the question posed by \citet{adebayo2020debugging,adebayo2022posthoc}, whose answer is
nuanced (detection tends to succeed for visible, known artifacts and to fail for non-visible or
unknown ones). Our aim is narrower and complementary. First, prior work largely treats the shortcut
as present or absent; we instead vary its \emph{strength} and ask how attribution-based detection
behaves along that dose--response axis. Second, we insist that a model's reliance be established
\emph{behaviorally} before any attribution is trusted, which turns out to be essential: we find an
attribution method that flags a marker the model does not actually use. Third, we compare three
standard methods---Grad-CAM \citep{selvaraju2017gradcam}, Integrated Gradients
\citep{sundararajan2017axiomatic}, and Occlusion Sensitivity \citep{zeiler2014visualizing}---under
identical conditions, and we are explicit about a spatial-resolution confound that makes one of the
three comparisons unfair as it stands.

Concretely, our contributions are: (i) a controlled, reproducible dose--response benchmark on a
medical image classifier, varying shortcut strength $\rho$ rather than toggling the shortcut on and
off; (ii) a detection protocol that couples above-chance attribution localization with a measured
behavioral reliance test, which we show guards against a real false-positive failure mode; and (iii)
a head-to-head comparison of the three attribution methods, reported with its limitations---including
the resolution confound affecting Grad-CAM---stated plainly rather than glossed over. All findings
are descriptive (three seeds, one binary task) and we make no claim of statistical significance or
of generalization beyond this setting.

% ===========================================================================
\section{Related Work}
\label{sec:related}

\paragraph{Shortcut learning and spurious correlations.}
\citet{geirhos2020shortcut} consolidate a range of generalization failures under the single notion
of shortcut learning and argue that benchmark accuracy can mask brittle, non-transferable decision
rules. The phenomenon is well documented in medical imaging specifically. \citet{zech2018variable}
train pneumonia classifiers across three hospital systems and show that models can recover the
source hospital almost perfectly and exploit site-level prevalence, so internal test performance
overstates real-world accuracy. \citet{degrave2021ai} find that COVID-19 chest-radiograph
classifiers rely on dataset-source confounders rather than pathology, and emphasize that external
validation alone does not guarantee a model uses clinically meaningful signal.
\citet{winkler2019association} demonstrate a shortcut with an unusually clear cause: gentian-violet
skin markings inflate a melanoma classifier's malignancy score for benign lesions, and cropping the
ink restores specificity. Our injected marker is a controlled analogue of this last, concrete case.

\paragraph{Attribution methods.}
We study three widely used post-hoc methods spanning the main families. Grad-CAM
\citep{selvaraju2017gradcam} weights final-layer convolutional activations by their class gradients
to produce a coarse, class-discriminative localization map. Integrated Gradients
\citep{sundararajan2017axiomatic} is a path-integral attribution defined by the Sensitivity and
Implementation-Invariance axioms, computed by integrating input gradients from a baseline to the
input. Occlusion Sensitivity \citep{zeiler2014visualizing} is a perturbation method that measures
the change in the target score as a window slides across and masks the image. These three differ in
mechanism (activation-gradient, path-integral, perturbation) and in spatial resolution, which is
central to interpreting our results.

\paragraph{Reliability and faithfulness of attributions.}
A substantial literature questions how far attributions can be trusted. \citet{adebayo2018sanity}
introduce randomization tests and show that some saliency methods are insensitive to model
parameters or to the data-generating process, behaving like edge detectors; notably, gradient-based
methods and Grad-CAM pass these checks while guided-backpropagation variants do not. In medical
imaging, \citet{arun2021assessing} report that eight saliency methods fail at least one
trustworthiness criterion on two public radiology datasets, and \citet{saporta2022benchmarking}
find that seven methods localize pathologies worse than a human benchmark, with Grad-CAM the
strongest of the group but still well below expert segmentation. These results motivate both our
chance-corrected localization metric and our refusal to treat a heat map as self-evidently
meaningful.

\paragraph{Detecting and correcting shortcuts with explanations.}
Explanations have been used to expose spurious reliance at the dataset level: \citet{lapuschkin2019unmasking}
cluster relevance maps to reveal ``Clever Hans'' behavior, such as a classifier keying on a source
tag in a subset of images. A separate line uses explanations to \emph{correct} shortcuts by
penalizing attribution in regions known to be irrelevant \citep{ross2017right,rieger2020interpretations};
critically, these methods presuppose that the shortcut has already been identified. The work closest
to ours asks whether post-hoc explanations can \emph{detect} an injected spurious signal in the
first place. \citet{adebayo2020debugging} find that attributions can help flag a visible
spurious-background bug, though a human study shows users lean on predictions rather than the maps;
\citet{adebayo2022posthoc} show that feature attribution, concept activation, and training-point
ranking are largely ineffective when the artifact is unknown at test time, especially for
non-visible artifacts, and that attribution can even indicate reliance where none exists. Our study
sits inside this nuance: we restrict attention to a visible, known shortcut, add a dose--response
axis that these works do not isolate, and reproduce---independently and in a medical setting---the
false-positive behavior that \citet{adebayo2022posthoc} warn about.

% ===========================================================================
\section{Background}
\label{sec:background}
Let $f$ be a trained classifier, $x$ an input image, and $f_c(x)$ the (pre-softmax) score for a
target class $c$. An attribution method produces a map $A_c(x)$ that assigns an importance value to
each input location.

\paragraph{Grad-CAM \citep{selvaraju2017gradcam}.} Let $A^k$ be the $k$-th feature map of the final
convolutional layer. Grad-CAM weights each feature map by the average gradient of the target score
with respect to it and keeps the positive evidence:
\begin{equation}
\alpha_k^c = \frac{1}{Z}\sum_{i,j}\frac{\partial f_c(x)}{\partial A^k_{ij}},
\qquad
L^c_{\text{Grad-CAM}} = \mathrm{ReLU}\!\left(\sum_k \alpha_k^c\, A^k\right),
\end{equation}
where $Z$ is the number of spatial locations. The map is produced at the \emph{resolution of the
feature map} and then upsampled to the image size; this resolution is central to our Grad-CAM
result (Section~\ref{sec:limitations}).

\paragraph{Integrated Gradients \citep{sundararajan2017axiomatic}.} For a baseline $x'$ (we use a
zero image), the attribution of input feature $i$ integrates the gradient along the straight path
from $x'$ to $x$:
\begin{equation}
\mathrm{IG}_i(x) = (x_i - x'_i)\int_{0}^{1}
\frac{\partial f_c\!\left(x' + \beta (x - x')\right)}{\partial x_i}\, d\beta,
\end{equation}
approximated by a Riemann sum over $50$ steps. IG satisfies the Sensitivity and
Implementation-Invariance axioms and yields a pixel-resolution map.

\paragraph{Occlusion Sensitivity \citep{zeiler2014visualizing}.} For a region $r$ (here a sliding
$4\times4$ window with stride $2$), the attribution is the drop in the target score when that region
is masked to a reference value:
\begin{equation}
A_c(r) = f_c(x) - f_c\!\left(x_{\setminus r}\right),
\end{equation}
where $x_{\setminus r}$ is $x$ with region $r$ occluded. Occlusion is model-agnostic and its
resolution is set by the window and stride rather than by any internal layer.

% ===========================================================================
\section{Research Questions and Hypotheses}
\label{sec:rq}
\begin{description}
  \item[RQ1 (Localization).] Do attribution maps place above-chance mass on the shortcut marker?
  \item[RQ2 (Detection).] Does a localization+behavior rule separate shortcut-reliant models from controls?
  \item[RQ3 (Consistency).] Do the three methods agree?
  \item[RQ4 (Robustness).] How does localization behave as shortcut strength $\rho$ varies?
\end{description}

% ===========================================================================
\section{Methodology}
\label{sec:method}

\paragraph{Dataset.} We use PneumoniaMNIST from MedMNIST v2 \citep{yang2023medmnist}, a binary
chest-X-ray task (normal vs.\ pneumonia) at $28\times28$ resolution with the official
train/validation/test split. Images are replicated to three channels and normalized with ImageNet
statistics (the backbone is ImageNet-pretrained; for the from-scratch model this is a fixed affine).

\paragraph{Controlled shortcut injection.} We stamp a solid $4\times4$-pixel high-intensity square
marker at a \emph{random} location (recorded per image as a ground-truth box $B$). Injection is
one-sided and controlled by a strength $\rho\in[0,1]$: a disease image (label $1$) receives a marker
with probability $\rho$; a normal image never does. Thus $\rho=0$ is a clean baseline (no markers)
and $\rho=1$ makes the marker a perfectly predictive cue. All injection decisions and positions are
drawn from seeded RNGs and are fully reproducible.

\paragraph{Models and training.} Each model is a classifier trained with AdamW (lr $10^{-3}$,
batch $128$, up to $10$ epochs, cross-entropy). The reported runs use an ImageNet-pretrained
ResNet-18 \citep{he2016resnet} (inputs upsampled to $64\times64$); a from-scratch SmallCNN is also
available. \textbf{Checkpoint selection uses a validation set drawn from the same biased
distribution as training}, so selection reflects a realistic in-distribution pipeline rather than
secretly rewarding the least shortcut-reliant model (which clean-validation selection would do).
We train $4~\text{strengths}\times 3~\text{seeds}=12$ models.

\paragraph{Test-time conditions.} On the held-out test split we evaluate: \emph{Clean} (no marker),
\emph{Concordant} (marker follows the training rule), and \emph{Conflict} (marker stamped on normal
images, breaking the correlation). Attribution is computed on the Conflict set for the model's
\emph{predicted} class.

\paragraph{Metrics.} Let $A(x)$ be an attribution map and $B$ the marker box.
\begin{align}
\mathrm{AM} &= \frac{\sum_{(i,j)\in B}|A(i,j)|}{\sum_{(i,j)}|A(i,j)|},
& \mathrm{AM}^{\star} &= \mathrm{AM} - \frac{|B|}{|x|},
\end{align}
where $|B|/|x|$ is the chance mass of an equal-area box. We also report the \emph{pointing-game}
hit rate (fraction of images whose attribution peak lies in $B$) and, per image, compare
$\mathrm{AM}$ to the $95$th percentile of $\mathrm{AM}$ over $500$ random equal-area boxes
(\emph{null-beat}). Behaviorally, we report clean accuracy, the accuracy gain
$\Delta\mathrm{Acc}=\mathrm{Acc}_{\text{concordant}}-\mathrm{Acc}_{\text{clean}}$, and the
\emph{adversarial flip rate} (AFR): among normal images classified correctly when clean, the
fraction that flip to ``pneumonia'' once the marker is added.

\paragraph{Detection rule.} A method \emph{detects} the shortcut in a model if it beats the random-box
null for a majority of Conflict images \textbf{and} the model is behaviorally reliant
($\mathrm{AFR}>0.3$). Requiring the behavioral condition guards against calling a model
shortcut-reliant purely because an attribution map responds to the visible marker.

% ===========================================================================
\section{Experimental Setup}
\label{sec:setup}
We use PyTorch \citep{paszke2019pytorch} and Captum \citep{kokhlikyan2020captum}. Integrated
Gradients uses a zero baseline and $50$ steps; Occlusion uses a $4\times4$ window with stride $2$
(matched to the marker). Attribution is evaluated on up to $100$ Conflict images per model. Seeds
$\{0,1,2\}$ control weights, data order, and marker placement; cuDNN is set deterministic. All code,
configurations, and per-model outputs are released; figures and tables below are generated directly
from \texttt{per\_model\_metrics.csv}.

% ===========================================================================
\section{Results}
\label{sec:results}

\paragraph{The injected shortcut induces measurable reliance (Table~\ref{tab:behavior}).}
At $\rho=0$ the models ignore the marker behaviorally (AFR $6.8\%$, $\Delta\mathrm{Acc}\approx 0$).
From $\rho=0.7$ onward the marker dominates behavior: AFR exceeds $90\%$, meaning the model flips
the large majority of healthy images to ``pneumonia'' when the marker is stamped on them. At
$\rho=1.0$ the shortcut also \emph{displaces} genuine learning: clean accuracy collapses from
$85\%$ to $48\%$ (with large seed variance), i.e.\ the model can barely classify unmarked images.

\begin{table}[t]
\centering
\caption{Behavioral reliance by shortcut strength $\rho$ (mean\,$\pm$\,SD over 3 seeds).}
\label{tab:behavior}
\begin{tabular}{lccc}
\toprule
$\rho$ & Clean acc.\ (\%) & $\Delta\mathrm{Acc}$ (\%) & AFR (\%) \\
\midrule
0.0 & $85.0\pm1.9$ & $-0.3\pm0.3$ & $6.8\pm6.1$ \\
0.7 & $90.7\pm1.3$ & $+0.6\pm0.2$ & $90.6\pm3.2$ \\
0.9 & $89.6\pm1.7$ & $+1.7\pm0.6$ & $91.2\pm6.9$ \\
1.0 & $48.0\pm15.2$ & $+51.7\pm15.8$ & $98.7\pm0.4$ \\
\bottomrule
\end{tabular}
\end{table}

\paragraph{IG and Occlusion localize the shortcut; Grad-CAM does not (Table~\ref{tab:xai}).}
Attribution mass above chance ($\mathrm{AM}^{\star}$) for Integrated Gradients and Occlusion rises
monotonically with $\rho$ and is well above zero for all shortcut-reliant models, with pointing-game
hit rates reaching $98\%$ and $90\%$ at $\rho=1.0$. Grad-CAM's $\mathrm{AM}^{\star}$ stays near zero
at every strength ($\leq 1.2\%$) and its peak almost never lands in the marker (hit rate $\approx
3\%$). Figure~\ref{fig:qual} shows a representative Conflict-set example illustrating this contrast.

\begin{table}[t]
\centering
\caption{Attribution localization by method and strength: $\mathrm{AM}^{\star}$ (mean\,$\pm$\,SD, \%),
pointing-game hit rate (\%), and models flagged by the detection rule (of 3 seeds).}
\label{tab:xai}
\begin{tabular}{lccc ccc ccc}
\toprule
 & \multicolumn{3}{c}{Grad-CAM} & \multicolumn{3}{c}{Integrated Gradients} & \multicolumn{3}{c}{Occlusion} \\
\cmidrule(lr){2-4}\cmidrule(lr){5-7}\cmidrule(lr){8-10}
$\rho$ & $\mathrm{AM}^{\star}$ & hit & det & $\mathrm{AM}^{\star}$ & hit & det & $\mathrm{AM}^{\star}$ & hit & det \\
\midrule
0.0 & $-0.2\pm0.1$ & 2  & 0/3 & $3.0\pm1.5$  & 22 & 0/3 & $0.1\pm0.3$  & 3  & 0/3 \\
0.7 & $1.0\pm0.4$  & 3  & 0/3 & $8.5\pm1.1$  & 75 & 3/3 & $7.8\pm2.0$  & 66 & 3/3 \\
0.9 & $0.9\pm0.5$  & 3  & 0/3 & $11.0\pm1.6$ & 84 & 3/3 & $13.7\pm2.6$ & 84 & 3/3 \\
1.0 & $1.2\pm0.3$  & 3  & 0/3 & $19.2\pm2.9$ & 98 & 3/3 & $17.2\pm6.0$ & 90 & 3/3 \\
\midrule
\multicolumn{1}{l}{\textbf{Total det.}} & & & \textbf{0/12} & & & \textbf{9/12} & & & \textbf{9/12} \\
\bottomrule
\end{tabular}
\end{table}

\begin{figure}[t]
\centering
\safefig[width=\linewidth]{qualitative_conflict.png}
\caption{Representative Conflict-set example at $\rho=1.0$ (the marker is stamped on a
\emph{normal} chest X-ray at a random location). Panels, left to right: input image; the
ground-truth marker box; and the Grad-CAM, Integrated Gradients, and Occlusion attribution maps for
the model's predicted class. Integrated Gradients and Occlusion concentrate on the marker, whereas
Grad-CAM is spatially diffuse and does not isolate it.}
\label{fig:qual}
\end{figure}

\begin{figure}[t]
\centering
\safefig{amstar_vs_rho.png}
\caption{Attribution mass above chance ($\mathrm{AM}^{\star}$) vs.\ shortcut strength $\rho$
(mean\,$\pm$\,SD over seeds). IG and Occlusion increase with $\rho$; Grad-CAM stays near zero.}
\label{fig:amstar}
\end{figure}

\paragraph{Localization tracks shortcut strength (RQ4).} Across the 12 models, $\mathrm{AM}^{\star}$
correlates with $\rho$ at Spearman $+0.95$ (IG) and $+0.89$ (Occlusion), and with the behavioral
displacement $\Delta\mathrm{Acc}$ at $+0.97$ and $+0.91$. Because AFR saturates at $\rho\geq0.7$,
within the reliant regime $\mathrm{AM}^{\star}$ tracks $\rho$ more than it tracks AFR
(Figure~\ref{fig:faith}). Grad-CAM shows only a weak, small-magnitude trend ($+0.61$ vs.\ $\rho$).

\begin{figure}[t]
\centering
\safefig{afr_vs_amstar.png}
\caption{Behavioral reliance (AFR) vs.\ localization ($\mathrm{AM}^{\star}$) across all 12 models.}
\label{fig:faith}
\end{figure}

\paragraph{A false-positive caution.} At $\rho=0$, Integrated Gradients already places above-chance
mass on the marker (null-beat $48\%$, hit $22\%$) even though the model does not rely on it
(AFR $6.8\%$). Attribution responding to a visible marker is therefore \emph{not} evidence of
reliance; the behavioral gate in our detection rule is what prevents a false alarm.

% ===========================================================================
\section{Discussion}
\label{sec:discussion}

We revisit the research questions in light of the results.

\paragraph{RQ1 (Localization) and RQ3 (Consistency).} Integrated Gradients and Occlusion
Sensitivity place clearly above-chance mass on the marker in every shortcut-reliant model, and they
agree with one another: both reach pointing-game hit rates near $90$--$98\%$ at $\rho=1.0$. Grad-CAM
does not---but, as discussed below, we treat this as a measurement artifact rather than a property
of the method. So two of the three methods are consistent localizers here, and the third is
inconclusive.

\paragraph{RQ4 (Robustness / dose--response).} The central positive finding is that localization is
not merely present or absent but \emph{graded}: $\mathrm{AM}^{\star}$ for IG and Occlusion rises
monotonically with shortcut strength (Spearman $+0.95$ and $+0.89$ against $\rho$;
Figure~\ref{fig:amstar}). This dose--response behavior is what a present/absent design cannot show,
and it suggests these attributions carry information about \emph{how much} a model leans on a cue,
not just whether it does. A caveat tempers the interpretation: because behavioral reliance (AFR)
saturates by $\rho=0.7$, within the reliant regime $\mathrm{AM}^{\star}$ tracks the \emph{training
correlation strength} more than it tracks the (already-maxed) behavioral reliance
(Figure~\ref{fig:faith}). Localization scales with the cue's statistical dominance, which is not the
same as a calibrated readout of causal dependence.

\paragraph{RQ2 (Detection) and the necessity of a behavioral gate.} Coupling localization with a
behavioral check cleanly separated the nine reliant models from the three controls. The behavioral
condition was not a formality: at $\rho=0$, Integrated Gradients already assigned above-chance mass
to the marker (null-beat $48\%$) even though the model ignored it (AFR $6.8\%$). An analyst reading
the heat map alone could have ``detected'' a shortcut the model does not use. This reproduces, in a
medical-imaging setting, the warning of \citet{adebayo2022posthoc} that attribution can indicate
reliance where none exists, and it argues concretely that a heat map should be corroborated by a
perturbation-based behavioral test before it is trusted.

\paragraph{Why IG and Occlusion succeed where Grad-CAM appears to fail.} The marker is a $4$-pixel
cue. IG operates at pixel resolution and Occlusion at the resolution of a $4$-pixel window, so both
can, in principle, resolve it; Grad-CAM operates at the resolution of the final convolutional
feature map, which on our ResNet-18 at $64\times64$ input is only $2\times2$. A map that coarse
cannot isolate a $4$-pixel region regardless of what the model learned. The Grad-CAM outcome is
therefore uninformative about the method's ability to detect shortcuts and should not be read as a
negative result; resolving it requires the ablation noted below. This also aligns with
\citet{saporta2022benchmarking}, who found Grad-CAM competitive among saliency methods but limited
for small targets.

\paragraph{Takeaway.} Within the easy regime of a visible, known shortcut, perturbation- and
path-integral attributions are useful detectors \emph{when paired with a behavioral test}, and their
signal scales with shortcut strength---consistent with the positive side of
\citet{adebayo2020debugging}. This is an encouraging but narrow result, and the limitations below
bound how far it should be carried.

% ===========================================================================
\section{Limitations}
\label{sec:limitations}
% These are load-bearing; keep them.
\begin{itemize}
  \item \textbf{Grad-CAM's failure is confounded by resolution.} On ResNet-18 with $64\times64$
  inputs the final convolutional feature map is only $2\times2$, which cannot resolve a $4$-pixel
  marker. We therefore \emph{do not} conclude that Grad-CAM is blind to shortcuts; a fair comparison
  requires a finer feature map (e.g.\ the SmallCNN variant, $14\times14$) or higher input resolution.
  \item \textbf{Small sample.} Three seeds and one binary task yield descriptive trends only; we make
  no claims of statistical significance or generalization across datasets/architectures.
  \item \textbf{Visible, known shortcut.} We study the easy regime; non-visible or unknown shortcuts
  (cf.\ \citet{adebayo2022posthoc}) are out of scope.
  \item \textbf{Method hyperparameters.} Occlusion window/stride and the IG baseline are fixed a
  priori; results may shift under other choices.
\end{itemize}

% ===========================================================================
\section{Conclusion}
\label{sec:conclusion}
We asked whether standard post-hoc attribution methods can detect a model's reliance on a shortcut,
in a controlled medical-imaging setting where the shortcut is known by construction. Across twelve
models spanning four shortcut strengths, Integrated Gradients and Occlusion Sensitivity localized
the injected marker in every shortcut-reliant model and in none of the controls, and their
localization scaled monotonically with shortcut strength---a dose--response signal that a
present/absent design cannot reveal. Crucially, attribution localization alone was not sufficient:
Integrated Gradients flagged the marker even in a control that did not use it, so a behavioral
reliance test was necessary to avoid false positives. Grad-CAM did not localize the marker, but we
attribute this to the coarse resolution of its feature map on our architecture rather than to a
fundamental limitation, and we identify a resolution-matched ablation as the key next step. These
findings are descriptive---three seeds on a single binary task with a visible, known
shortcut---and should be read as evidence about that regime rather than a general guarantee.
The broader lesson is practical: a heat map that highlights a suspicious region is a hypothesis, not
a verdict, and should be confirmed behaviorally before a model is trusted or rejected.

\bibliographystyle{plainnat}
\bibliography{references}

\end{document}
